% 12 · EVoC: embedding, density, and persistence in one object
\section{EVoC: ONE FIT, THREE STAGES}
\label{sec:evoc}

\deckslides{slides 43--45: the fusion, the pipeline, and why C-space.}

The three methods so far are composed by hand: t-SNE or UMAP produces an
embedding, then HDBSCAN$^*$ (or a persistence-driven variant of it) produces
labels, then the user decides how the two are coupled. EVoC
\citep{EVoC} removes the seam. It takes the graph, embeds it, clusters it, and
chooses the clustering itself, in a single fit, and it is the method our
pipeline uses for its production runs.

Read the deck's pipeline figure first: the point of the method is that the
three stages share one objective rather than being tuned against each other.

\begin{figure}[htbp]\centering
  \fig[0.92]{evoc_pipeline.png}
  \caption{The EVoC pipeline. From left to right: a nearest-neighbour graph is
  built from the raw data with a native metric, the graph is embedded into a
  low-dimensional space by a graph-embedding objective, and the embedded points
  are clustered with a density-based method whose candidate clusterings are
  scored by persistence. The output is labels, not a picture.}
  \label{fig:evocpipeline}
\end{figure}

\subsection{Stage 1: the graph, built on the metric you meant}

EVoC's first stage is a $k$-nearest-neighbour graph (\S\ref{sec:knn})
constructed directly from the input representation, with a configurable
metric. This matters in our setting for a specific reason: it is the one arm
that can use the \emph{raw} L2-normalised abundance matrix and its native
cosine geometry, without an intervening Euclidean-only embedding. Since
\S\ref{sec:data} normalises every row to unit length, Euclidean and cosine
distances are monotonically related on our matrix, and EVoC's use of cosine is
a deliberate choice rather than an accident: in a compositional space, only
the ratios between elements are meaningful, and a vector's length is an
artefact of the abundance sum.

\begin{figure}[htbp]\centering
  \begin{tabular}{@{}c@{\hskip6pt}c@{}}
    \fig[0.44]{evoc_step1.png} &
    \fig[0.44]{evoc_step2.png} \\[3pt]
    \fig[0.44]{evoc_step3.png} &
    \fig[0.44]{evoc_step4.png}
  \end{tabular}
  \caption{The four stages as the lecture builds them: \panel{a} the
  $k$-nearest-neighbour graph of the data; \panel{b} an embedding of the graph
  nodes (the graph's structure is preserved, not its geometry); \panel{c}
  density-based clustering of the embedded points; \panel{d} the persistence
  score used to pick which clustering layer to return.}
  \label{fig:evocsteps}
\end{figure}

\subsection{Stage 2: an embedding for clustering, not for looking at}

The embedding stage optimises a graph-embedding objective: adjacent nodes
close, non-adjacent nodes apart, with the dimension chosen by the algorithm
rather than by the user. It is worth being explicit about what this is
\emph{not}. It is not a visualisation; there is no meaningful scatter plot at
the end of it, and we never draw one. Its purpose is to place the graph in a
space where a density-based clusterer can separate its communities, which is a
different objective from t-SNE's or UMAP's two-dimensional layout problem.

This is also where the most common confusion in the literature arises. When a
paper says ``UMAP followed by HDBSCAN$^*$'', there are two algorithms and
therefore two independent sources of error; when a paper says ``EVoC'', the
cluster structure is chosen inside the fit, by persistence, from the same graph
that the embedding was built on. The distinction is not about quality but about
where the arbitrary decisions live, and it is why our benchmark reports all
three, since each one answers a different question about the same data.

\subsection{Stage 3: layers, persistence, and no \texttt{min\_cluster\_size}}

The clustering stage produces not one answer but a family of them, indexed by
scale. Each candidate layer is scored by the total persistence of the clusters
it contains (the PLSCAN criterion of \S\ref{sec:plscan}) and the best-scoring
layer's labels are returned. Consequences that matter in practice:

\begin{itemize}
\item \textbf{No size floor to set.} The user is never asked for
  \texttt{min\_cluster\_size}, which is the single knob that most changed
  results in our DR17 sweeps (\S\ref{sec:benchmark}: raising it from 5 to 10
  moved t-SNE recall from $0.087$ to $0.404$). EVoC replaces it with an
  internal persistence criterion, which is more stable across datasets but
  less inspectable: you cannot see which layers were considered.
\item \textbf{The number of clusters is an output.} On the field runs, EVoC
returns a few dozen groups on average, none of them large: one of the reasons
its recall is high and its precision catastrophic
(\textbf{Table~\ref{tab:field}}).
\item \textbf{The returned layer is a model-selection decision.} Persistence
  is a principled criterion, not a guarantee: if the field and the clusters
  are not separated in the graph's geometry, the best-scoring layer is the one
  that splits the field most persistently.
\end{itemize}

\subsection{What EVoC does for us, and what it does not}

Positive first. EVoC is the only arm that needs no post-hoc clustering choice,
it works directly on the cosine geometry that our data actually have, it is
the cheapest of the three to run on large samples, and on cluster-only
separation it is competitive: homogeneity $0.420 \pm 0.043$ on abundances,
$0.733 \pm 0.008$ on the 64-dimensional PCA latent, both measured over seven
seeds on the same 982 stars as the other arms.

Now the part that took us a while to accept. Two results in the literature
around this method, and two of ours, all point the same way:

\begin{issues}[WHAT WE HAVE TO SAY AGAINST OUR OWN FAVOURITE] \textbf{The seed
spread is larger than several of the gaps.} EVoC's homogeneity on abundances is
$0.420 \pm 0.043$ over seven seeds: a spread that overlaps with the t-SNE and
UMAP arms. In the five-cluster head-to-head the spread is $\pm0.08$--$0.11$.
Any ranking that depends on differences of that size is not a ranking.

\textbf{The self-supervision claim is a t-SNE/UMAP result, not an EVoC result.}
On the corrected same-population table, the masked autoencoder scores $0.70 \pm
0.08$ against $0.64 \pm 0.08$ for abundances on EVoC (overlapping within
errors) while the \emph{supervised} CNN wins EVoC outright at $0.80 \pm 0.03$
on the four-cluster sample. We withdrew the sentence that said the
abundance-free latent wins everywhere.

\textbf{High recall at near-zero precision.} On the all-sky field-retrieval
task, EVoC recovers $0.557$ of the true members at a precision of $0.0033$:
the groups it returns are mostly field stars, by two orders of magnitude. It
is the same failure that UMAP shows in that table, and it is the reason the
pipeline's field work runs a rejection stage after the clustering rather than
trusting the labels.
\end{issues}

\begin{issues}[EXERCISES]
\begin{enumerate}
\item EVoC clusters an embedding built from the graph, and HDBSCAN$^*$
  clusters the embedding built by t-SNE or UMAP. Explain, in three sentences
  each, why the first is more self-consistent and why the second is easier to
  diagnose when it fails.
\item On the member matrix, run EVoC with the seed varied over the same seven
  values used in this workbook, and report homogeneity as mean $\pm$ standard
  deviation. Now do the same for t-SNE with a fixed seed but nine row orders
  (\S\ref{sec:validation}). Which source of variability is larger, and what
  does that say about which number belongs in a paper?
\item Take the EVoC labels on the all-sky field run and compute, for each
  returned group, the fraction of its members that are known field stars.
  Plot that distribution. How many groups would you have to ignore before the
  remaining ones have precision above 0.5?
\item The persistence criterion selects a clustering layer without user
  input. Construct a two-dimensional dataset with a dense uniform background
  and one small dense cluster, and report which layer EVoC selects. Does the
  criterion find the cluster, the background, or something else?
\end{enumerate}
\end{issues}
